\documentclass[../../../include/open-logic-section]{subfiles}

\begin{document}
\olfileid[id]{sth}{card-arithmetic}{simp}

\olsection{Menyederhanakan Penjumlahan dan Perkalian}

Ternyata penjumlahan dan perkalian kardinal transfinit \emph{sangat} mudah.
Hal ini mengikuti dari fakta bahwa kardinal merupakan ordinal (tertentu),
sehingga terurut baik dan dapat dimanipulasi dengan cara tertentu. Namun,
menunjukkan hal ini \emph{tidak} begitu mudah. Sebagai permulaan, kita
memerlukan suatu definisi yang agak rumit:

\begin{defn}
Kita mendefinisikan suatu \emph{pengurutan kanonik}, $\canonord$, pada
pasangan ordinal dengan menetapkan bahwa
$\tuple{\alpha_1, \alpha_2} \canonord
\tuple{\beta_1, \beta_2}$ jika dan hanya jika salah satu dari hal-hal berikut
berlaku:
\begin{enumerate}
	\item $\max(\alpha_1, \alpha_2) < \max(\beta_1, \beta_2)$; atau
	\item $\max(\alpha_1, \alpha_2) = \max(\beta_1, \beta_2)$ dan
	$\alpha_1 < \beta_1$; atau
	\item $\max(\alpha_1, \alpha_2) = \max(\beta_1, \beta_2)$ dan
	$\alpha_1 = \beta_1$ dan $\alpha_2 < \beta_2$
\end{enumerate}
\end{defn}

\begin{lem}
$\tuple{\alpha \times \alpha, \canonord}$ merupakan suatu pengurutan baik
untuk sebarang ordinal $\alpha$.
\end{lem}

\begin{proof}
Jelas bahwa $\canonord$ terhubung pada $\alpha \times \alpha$. Andaikan
baik $\tuple{\alpha_1, \alpha_2}$ maupun $\tuple{\beta_1, \beta_2}$ tidak
$\canonord$-lebih kecil daripada yang lain. Maka
$\max(\alpha_1, \alpha_2) = \max(\beta_1, \beta_2)$,
$\alpha_1 = \beta_1$, dan $\alpha_2 = \beta_2$, sehingga
$\tuple{\alpha_1, \alpha_2} = \tuple{\beta_1,  \beta_2}$.

Untuk menunjukkan pengurutan baik, misalkan
$X \subseteq \alpha\times\alpha$ tak kosong. Karena $\alpha$ suatu ordinal,
terdapat suatu $\delta$ yang merupakan anggota terkecil dari
$\Setabs{\max(\gamma_1, \gamma_2)}{\tuple{\gamma_1, \gamma_2} \in X}$.
Sekarang, dari
$\Setabs{\tuple{\gamma_1,\gamma_2} \in X}{\max(\gamma_1, \gamma_2) =
\delta}$, singkirkan semua pasangan selain yang mempunyai koordinat pertama
terkecil; dari pasangan-pasangan yang tersisa, pasangan dengan koordinat
kedua terkecil merupakan anggota $\canonord$-terkecil dari $X$.
\end{proof}
\noindent
Sekarang untuk suatu pengamatan kecil dan sederhana:

\begin{prop}\ollabel{simplecardproduct}
Jika $\cardeq{\alpha}{\beta}$, maka
$\cardeq{\alpha \times \alpha}{\beta \times \beta}$.
\end{prop}

\begin{proof}
Cukup misalkan $f \colon \alpha \to \beta$ menginduksi pemetaan
$\tuple{\gamma_1, \gamma_2} \mapsto
\tuple{f(\gamma_1), f(\gamma_2)}$.
\end{proof}
\noindent
Sekarang kita akan menggunakan semua ini untuk membuktikan suatu lema
penting:
\begin{lem}\ollabel{alphatimesalpha}
$\cardeq{\alpha}{\alpha \times \alpha}$, untuk sebarang ordinal tak hingga
$\alpha$
\end{lem}

\begin{proof}
Untuk memperoleh kontradiksi, misalkan $\alpha$ ordinal tak hingga terkecil
yang tidak memenuhi pernyataan tersebut.
\olref[sfr][siz][zigzag]{natsquaredenumerable} menunjukkan bahwa
$\cardeq{\omega}{\omega\times\omega}$, sehingga $\omega \in \alpha$.
Selain itu, $\alpha$ merupakan suatu kardinal. Untuk memperoleh kontradiksi,
andaikan bukan demikian; maka $\card{\alpha} \in \alpha$, sehingga
$\cardeq{\card{\alpha}}{\card{\alpha} \times \card{\alpha}}$ berdasarkan
hipotesis minimalitas; dan $\cardeq{\card{\alpha}}{\alpha}$ berdasarkan
definisi; sehingga $\cardeq{\alpha}{\alpha\times\alpha}$ berdasarkan
\olref{simplecardproduct}.

Sekarang, untuk setiap
$\tuple{\gamma_1, \gamma_2} \in \alpha \times \alpha$, perhatikan segmen:
\begin{align*}
	\text{Seg}(\gamma_1, \gamma_2) &= \Setabs{\tuple{\delta_1, \delta_2} \in \alpha \times \alpha}{\tuple{\delta_1, \delta_2} \canonord \tuple{\gamma_1, \gamma_2}}
\end{align*}
Dengan menetapkan $\gamma = \max(\gamma_1, \gamma_2)$, perhatikan bahwa
$\tuple{\gamma_1, \gamma_2} \canonord
\tuple{\gamma+1, \gamma + 1}$. Jadi, ketika $\gamma$ tak hingga, perhatikan:
\begin{align*}
	\text{Seg}(\gamma_1, \gamma_2) & 
	\precsim ((\gamma \ordplus 1)\ordtimes (\gamma \ordplus 1))\\
	&\approx (\gamma \ordtimes \gamma)
	\text{, berdasarkan \olref[ord-arithmetic][using-addition]{ordinfinitycharacter} dan 
	\olref{simplecardproduct}}\\
	&\approx \gamma \text{, berdasarkan hipotesis induksi}\\
	& \prec \alpha\text{, sebab $\alpha$ merupakan suatu kardinal}
\end{align*}
Jadi $\ordtype{\alpha\times \alpha, \canonord} \leq \alpha$, dan dengan
demikian $\cardle{\alpha \times \alpha}{\alpha}$. Karena tentu saja
$\cardle{\alpha}{\alpha \times \alpha}$, hasilnya mengikuti dari
Schr\"oder--Bernstein.
\end{proof}

Akhirnya, kita memperoleh hasil penyederhanaan kita:

\begin{thm}\ollabel{cardplustimesmax}
Jika $\cardfont{a}, \cardfont{b}$ merupakan kardinal tak hingga, maka:
\[
	\cardfont{a}
\cardtimes \cardfont{b} = \cardfont{a} \cardplus \cardfont{b} =
\text{max}(\cardfont{a}, \cardfont{b}).
\]
\end{thm}

\begin{proof}
Tanpa mengurangi keumuman, andaikan
$\cardfont{a} = \max(\cardfont{a}, \cardfont{b})$. Kemudian, dengan
menggunakan \olref{alphatimesalpha},
$\cardfont{a}\cardtimes\cardfont{a} = \cardfont{a} \leq \cardfont{a}
\cardplus \cardfont{b} \leq \cardfont{a} \cardplus \cardfont{a} \leq
\cardfont{a} \cardtimes \cardfont{a}$. \end{proof}\noindent Serupa dengan
itu, jika $\cardfont{a}$ tak hingga, suatu gabungan berukuran
$\cardfont{a}$ dari himpunan-himpunan yang masing-masing berukuran
$\leq\cardfont{a}$ mempunyai ukuran $\leq\cardfont{a}$:

\begin{prop}\ollabel{kappaunionkappasize}
Misalkan $\cardfont{a}$ suatu kardinal tak hingga. Untuk setiap ordinal
$\beta \in \cardfont{a}$, misalkan $X_\beta$ suatu himpunan dengan
$\card{X_\beta} \leq \cardfont{a}$. Maka
$\card{\bigcup_{\beta \in \cardfont{a}} X_\beta} \leq \cardfont{a}$.
\end{prop}

\begin{proof}
Untuk setiap $\beta \in \cardfont{a}$, tetapkan !!a{injection}
$f_\beta \colon X_\beta \to \cardfont{a}$.\footnote{Bagaimana semuanya
``ditetapkan''? Lihat \olref[sth][choice][countablechoice]{sec}.} Definisikan
!!a{injection} $g \colon
\bigcup_{\beta \in \cardfont{a}} X_\beta \to \cardfont{a} \times
\cardfont{a}$ dengan $g(v) = \tuple{\beta, f_\beta(v)}$, di mana
$v \in X_\beta$ dan $v \notin X_\gamma$ untuk setiap
$\gamma \in \beta$. Sekarang
$\bigcup_{\beta \in \cardfont{a}} X_\beta \preceq \cardfont{a} \times
\cardfont{a} \approx \cardfont{a}$ berdasarkan
\olref{cardplustimesmax}.
\end{proof}

\end{document}
